% !TEX root = ../main.tex
% !TEX program = lualatex
\documentclass[../main.tex]{subfiles}
\IfSubfilesClassLoaded{\externaldocument{\subfix{../build/main}}}{}
% ====================
% File: 17I_CM_and_Technical_Reviews.tex
% ====================
\begin{document}
\IfSubfilesClassLoaded{\chapter{Configuration Management and Technical Reviews}}{}
% === SE Analysis & Control: CM & Technical Reviews (EDO 3.5.3) ===
\section{SE Analysis \& Control: Configuration Management \& Technical Reviews}

\subsection{Summary}
\begin{description}
  \item[\textbf{CM purpose.}] Configuration management establishes and maintains consistency between requirements, the product, and configuration information across the life cycle~\autocite{edo-3-5-3-cm-tech-reviews-2025}.
  \item[\textbf{CI control.}] Designating \acp{ci} enables separate control, documentation, reviews, and qualification testing~\autocite{edo-3-5-3-cm-tech-reviews-2025}.
  \item[\textbf{Baselines.}] Functional, allocated, and product baselines connect specifications to technical reviews and change control~\autocite{edo-3-5-3-cm-tech-reviews-2025}.
  \item[\textbf{Interfaces and data.}] Interface control and technical data management preserve configuration discipline across government and contractor boundaries~\autocite{edo-3-5-3-cm-tech-reviews-2025}.
  \item[\textbf{Technical reviews.}] Reviews assess design maturity, readiness, compliance, and risk at key points in the acquisition life cycle~\autocite{edo-3-5-3-cm-tech-reviews-2025}.
\end{description}

\subsection{Configuration Management}
\begin{description}
  \item[\textbf{Definition.}] A technical and management process that maintains consistency between product requirements, the product, and configuration information across baselines~\autocite{edo-3-5-3-cm-tech-reviews-2025}.
  \item[\textbf{Implementation.}] Technical and administrative direction to identify and document \acp{ci}, control changes, and record change status to provide an audit trail~\autocite{edo-3-5-3-cm-tech-reviews-2025}.
\end{description}

\subsection{Why CM Matters}
\begin{itemize}
  \item Provides measurable performance parameters and a common basis for acquisition and use~\autocite{edo-3-5-3-cm-tech-reviews-2025}.
  \item Ensures correct, current information for change decisions and minimizes avoidable cost~\autocite{edo-3-5-3-cm-tech-reviews-2025}.
  \item Correlates products with requirements, design, and product information to avoid guesswork~\autocite{edo-3-5-3-cm-tech-reviews-2025}.
  \item Verifies configuration against required attributes and increases confidence in product information~\autocite{edo-3-5-3-cm-tech-reviews-2025}.
\end{itemize}

\subsection{Configuration Items}
\begin{description}
  \item[\textbf{CI definition.}] A hardware/software aggregation treated as a single entity for configuration management and uniquely identified at a reference point~\autocite{edo-3-5-3-cm-tech-reviews-2025}.
  \item[\textbf{Designation impacts.}] Separate control, documentation, formal approval of changes, individual design reviews, and qualification tests~\autocite{edo-3-5-3-cm-tech-reviews-2025}.
\end{description}

\subsection{Considerations for Designating CIs}
\begin{itemize}
  \item Qualification level, integration of tested/untested components, and interfaces with other controlled \acp{ci}~\autocite{edo-3-5-3-cm-tech-reviews-2025}.
  \item Need to know exact configuration during the life cycle, high-risk or safety concerns, and new technology~\autocite{edo-3-5-3-cm-tech-reviews-2025}.
  \item Each development contract should designate at least one \ac{ci}; cost alone is not a valid consideration~\autocite{edo-3-5-3-cm-tech-reviews-2025}.
\end{itemize}

\subsection{CM Functional Components}
\begin{description}
  \item[\textbf{Identification.}] Identify and document functional and physical characteristics of \acp{ci}~\autocite{edo-3-5-3-cm-tech-reviews-2025}.
  \item[\textbf{Status accounting.}] Record and report all configuration changes~\autocite{edo-3-5-3-cm-tech-reviews-2025}.
  \item[\textbf{Control.}] Manage \acp{ci} and documentation throughout the life cycle~\autocite{edo-3-5-3-cm-tech-reviews-2025}.
  \item[\textbf{Verification and audits.}] Verify conformance to requirements, specifications, drawings, and interface documents~\autocite{edo-3-5-3-cm-tech-reviews-2025}.
\end{description}

\subsection{IPPD and CM}
\begin{description}
  \item[\textbf{Baseline linkage.}] CM ties requirements and specifications to baselines and supports \ac{ippd} across disciplines~\autocite{edo-3-5-3-cm-tech-reviews-2025}.
  \item[\textbf{Cross-discipline impact.}] Configuration changes affect test and evaluation, manufacturing, logistics, software development, and technical data management~\autocite{edo-3-5-3-cm-tech-reviews-2025}.
\end{description}

\subsection{Commercial Item Impacts}
\begin{description}
  \item[\textbf{Government control limits.}] Design details and full \ac{tdp} may not be available, complicating manuals and sustainment~\autocite{edo-3-5-3-cm-tech-reviews-2025}.
  \item[\textbf{Interfaces and support.}] Interfaces must be carefully defined, and support strategies must address commercial repair and updates~\autocite{edo-3-5-3-cm-tech-reviews-2025}.
  \item[\textbf{Test strategy.}] Tests must ensure performance at both item and integrated levels~\autocite{edo-3-5-3-cm-tech-reviews-2025}.
\end{description}

\subsection{Configuration Control}
\begin{description}
  \item[\textbf{Engineering change.}] Modifies, adds, deletes, or supersedes original parts~\autocite{edo-3-5-3-cm-tech-reviews-2025}.
  \item[\textbf{Deviation.}] Written authorization before manufacture to depart from requirements for specific units or time~\autocite{edo-3-5-3-cm-tech-reviews-2025}.
  \item[\textbf{Waiver.}] Written authorization during or after manufacture to accept a nonconforming item as-is or after rework~\autocite{edo-3-5-3-cm-tech-reviews-2025}.
  \item[\textbf{Approved configuration.}] Approved baseline plus approved changes~\autocite{edo-3-5-3-cm-tech-reviews-2025}.
\end{description}

\subsection{Configuration Baselines}
\begin{description}
  \item[\textbf{Functional baseline.}] Defines approved system-level functional, performance, interoperability, interface, and verification requirements; the program manager normally establishes Government control at \ac{sfr}~\autocite{edo-3-5-3-cm-tech-reviews-2025,DAU-Functional-Baseline}.
  \item[\textbf{Allocated baseline.}] Defines the allocated performance, interface, derived, design-restraint, and verification requirements for the \acp{ci} that make up the system; it is normally established at \ac{pdr}~\autocite{edo-3-5-3-cm-tech-reviews-2025,DAU-Allocated-Baseline}.
  \item[\textbf{Product baseline.}] Defines the detailed build-to or code-to functional and physical characteristics of \acp{ci}; the initial product baseline is normally established at \ac{cdr} and finalized through \ac{pca}~\autocite{edo-3-5-3-cm-tech-reviews-2025,DAU-Product-Baseline,DAU-Critical-Design-Review}.
\end{description}

\subsection{Specifications, CM Planning, and Baselines}
\begin{description}
  \item[\textbf{System specification.}] Establishes mission performance requirements and interfaces tied to the functional baseline~\autocite{edo-3-5-3-cm-tech-reviews-2025}.
  \item[\textbf{Performance specification.}] Defines form, fit, and function of \acp{ci} and supports the allocated baseline~\autocite{edo-3-5-3-cm-tech-reviews-2025}.
  \item[\textbf{Detail/process/material specifications.}] Define build-to requirements and fabrication processes that support the product baseline~\autocite{edo-3-5-3-cm-tech-reviews-2025}.
\end{description}

\subsection{Interface Control}
\begin{description}
  \item[\textbf{Interfaces.}] Boundaries between \acp{ci} or between contractor and \ac{gfe} requiring mechanical, electrical, operational, or software alignment~\autocite{edo-3-5-3-cm-tech-reviews-2025}.
  \item[\textbf{ICWG.}] The \ac{icwg} identifies, documents, and controls interface characteristics and produces interface control documents~\autocite{edo-3-5-3-cm-tech-reviews-2025}.
  \item[\textbf{Government role.}] Government sets interface standards; contractors design to meet interface specifications~\autocite{edo-3-5-3-cm-tech-reviews-2025}.
\end{description}

\subsection{Technical Data Management}
\begin{description}
  \item[\textbf{Purpose.}] Processes to plan, acquire, manage, and use technical data to support the system life cycle~\autocite{edo-3-5-3-cm-tech-reviews-2025}.
  \item[\textbf{Data requirements.}] Documented in the \ac{cdrl}; only minimum required data should be procured in contractor format~\autocite{edo-3-5-3-cm-tech-reviews-2025}.
  \item[\textbf{TDP support.}] \acp{tdp} support technical reviews, including \ac{fca} and \ac{pca}~\autocite{edo-3-5-3-cm-tech-reviews-2025}.
\end{description}

\subsection{Configuration Management Strategy}
\begin{description}
  \item[\textbf{Government-controlled.}] Functional baseline, system-level specs, external interfaces, and system or major subsystem reviews~\autocite{edo-3-5-3-cm-tech-reviews-2025}.
  \item[\textbf{Contractor-controlled.}] Allocated and product baselines, item performance/detail specs, internal interfaces, and component/vendor reviews~\autocite{edo-3-5-3-cm-tech-reviews-2025}.
  \item[\textbf{Benefits.}] Promotes design flexibility, reduces government burden for low-level \acp{ecp}, and supports smaller program offices~\autocite{edo-3-5-3-cm-tech-reviews-2025}.
\end{description}

\subsection{Technical Reviews Purpose}
\begin{itemize}
  \item Check design maturity and readiness for the next development level~\autocite{edo-3-5-3-cm-tech-reviews-2025}.
  \item Assess compliance with requirements and review technical risk~\autocite{edo-3-5-3-cm-tech-reviews-2025}.
  \item Provide visibility into contractor implementation of the statement of work~\autocite{edo-3-5-3-cm-tech-reviews-2025}.
\end{itemize}

\subsection{MCA Technical Reviews by Phase}
\ac{mca} technical reviews are event-driven maturity points, not calendar ceremonies. Unless waived through the \ac{sep} approval process, \ac{dodi}~5000.88 requires the \ac{pm} to conduct system-level reviews or equivalents for requirements/functionality, preliminary design, critical design, verification/configuration audit, production readiness, and physical configuration audit. DAU's technical-baseline maturation table separates \ac{srr} as mutual requirements understanding from \ac{sfr} as functional-baseline approval; \ac{pdr} and \ac{cdr} then anchor the allocated and initial product baselines~\autocite{DoDI5000-88,DoD-EngineeringDefenseSystemsGuidebook-2024,DAU-Technical-Baselines,DAU-System-Requirements-Review,DAU-System-Functional-Review,DAU-Preliminary-Design-Review,DAU-Critical-Design-Review}.

\begin{longtblr}[
    label   = {tab:mca-technical-reviews},
    caption = {Major Capability Acquisition technical reviews by phase},
    entry   = {MCA Technical Reviews by Phase},
    remark{\tSource} = {\tabCite{DoDI5000-88,DoD-EngineeringDefenseSystemsGuidebook-2024,DAU-Technical-Baselines,DoDI5000-89,DoDI5000-98,OPNAVINST-5450-332C}}
  ]
  {colspec = {@{} l X[1.15] X[1.25] X[3.2] @{}}}
  \toprule
  {Review} & {Typical MCA phase}                                  & {Primary baseline/readiness focus}                    & {Board-prep purpose}                                                                                                                                                            \\
  \midrule
  ASR      & MSA                                                  & Preferred materiel solution                           & Confirms the preferred materiel solution is technically feasible, affordable, and understood well enough to support a draft performance specification.                          \\
  SRR      & TMRR                                                 & System requirements                                   & Confirms the developer and government share a clear, testable, traceable understanding of system performance requirements.                                                      \\
  SFR      & TMRR                                                 & Functional baseline                                   & Confirms the functional architecture satisfies performance requirements and supports preliminary design with acceptable risk.                                                   \\
  PDR      & TMRR, normally before Milestone B                    & Allocated baseline                                    & Confirms the preliminary design and architecture are mature enough to proceed to detailed design and support the Milestone B decision.                                          \\
  CDR      & EMD                                                  & Product baseline                                      & Confirms detailed design stability, expected performance, affordability posture, and readiness to start coding, fabrication, integration, and test.                             \\
  TRR      & EMD, before major test events                        & Test readiness                                        & Confirms the test article, procedures, instrumentation, resources, environment, safety posture, and success criteria are ready for the planned test.                            \\
  SVR/FCA  & EMD                                                  & Verification against functional or allocated baseline & Verifies the system or configuration item satisfies approved requirements and that discrepancies are understood before production or operational test decisions.                \\
  PRR      & EMD before LRIP; repeated in P\&D when needed        & Production readiness                                  & Confirms producibility, quality, manufacturing planning, supply chain, facilities, tooling, and production risk are acceptable for LRIP or FRP.                                 \\
  OTRR     & P\&D, before OPEVAL/IOT\&E or follow-on OT\&E         & Operational test readiness                            & Certifies that the production-representative configuration, reliability, maintainability, supportability, hazards, resources, and test planning are ready for operational test. \\
  PCA      & P\&D, after successful OT\&E and before FRP decision & Final product baseline                                & Audits the as-built production-representative item against its technical documentation, resolves discrepancies, and establishes the final product baseline.                     \\
  \bottomrule
\end{longtblr}

\Fig{\ref{fig:mca_schedule}} shows when the major technical reviews occur.

\note{For a board answer, lead with the baseline progression: \ac{srr} reviews the system requirements, \ac{sfr} supports the functional baseline, \ac{pdr} supports the allocated baseline, \ac{cdr} supports the initial product baseline, \ac{svr}/\ac{fca} verify configuration, \ac{prr} checks production readiness, \ac{otrr} checks operational test readiness, and \ac{pca} supports the final as-built product baseline.}

\subsection{System-Level Technical Review Board Details}
For oral boards, answer each review with four parts: why the review exists, who must be in the room, when it occurs, and what decision or baseline it supports. The \ac{pm} owns the review plan through the \ac{sep}; the chief engineer or lead systems engineer normally drives technical closure; \acp{ipt}, contractors, users, logisticians, test, manufacturing, cybersecurity, configuration-management, cost, contracting, and independent subject-matter experts provide the evidence. For \ac{acat}~ID design reviews, the \ac{pm} invites \ac{usdre} representatives, and independent review teams advise the \ac{cae} without being controlled by the \ac{pmo}~\autocite{DoDI5000-88,DoD-EngineeringDefenseSystemsGuidebook-2024}.

\begin{longtblr}[
    label   = {tab:technical-review-board-details},
    caption = {System-level technical review board details},
    entry   = {System-Level Technical Review Board Details},
    remark{\tSource} = {\tabCite{DoDI5000-88,DoD-EngineeringDefenseSystemsGuidebook-2024,DAU-Technical-Baselines,DoDI5000-89,DoDI5000-98,OPNAVINST-5450-332C}}
  ]
  {colspec = {@{} l X[1.35] X[1.55] X[1.9] @{}}}
  \toprule
  {Review} & {Major players} & {When it normally occurs} & {Milestone, event, or board cue} \\
  \midrule
  ASR & PM, systems engineer, sponsor/user, AoA and cost teams, technical authority, sustainment, test, and independent reviewers & Late MSA, after AoA work identifies a preferred materiel solution & Supports Milestone A/TMRR planning by confirming the preferred solution is feasible, affordable, and mature enough to guide draft performance requirements. \\
  SRR & PM, systems engineer, requirements sponsor/user, contractor, test, logistics, cybersecurity, configuration management, and independent reviewers & TMRR, after system requirements and the system performance specification are mature enough to review & Confirms requirements are clear, testable, traceable, and mutually understood; do not call it the normal baseline-setting event. \\
  SFR & PM, systems engineer, architecture/function leads, contractor, users, test, logistics, cybersecurity, and configuration management & TMRR, after functional analysis and system architecture mature beyond SRR & Supports the functional baseline and approval of the system performance specification for configuration control. \\
  PDR & PM, systems engineer, contractor design leads, manufacturing, logistics, test, cost, risk, independent reviewers, and for ACAT ID OUSD(R\&E) & Late TMRR, normally before Milestone B unless waived & Supports the Milestone B decision and EMD RFP/contracting posture by establishing the allocated baseline and confirming readiness for detailed design. \\
  CDR & PM, chief engineer or lead systems engineer, contractor design leads, manufacturing, logistics, test, configuration management, independent reviewers, and for ACAT ID OUSD(R\&E) & EMD, after detailed design and build-to/code-to data are mature & Establishes the initial product baseline and authorizes fabrication, coding, integration, and test with acceptable residual risk. \\
  TRR & PM, chief developmental tester, lead DT organization, systems engineer, safety, range/lab, instrumentation, contractor, and certification authorities & Before major DT, qualification, cybersecurity, interoperability, or live test events & Gates test execution by confirming the article, environment, procedures, resources, safety posture, instrumentation, and success criteria are ready. \\
  SVR/FCA & PM, systems engineer, test, configuration management, contractor, certification authorities, and independent technical reviewers & EMD, after integration and verification evidence are mature & Verifies the system or configuration item satisfies approved functional and allocated requirements and identifies discrepancies before production or operational-test decisions. \\
  PRR & PM, systems engineer, manufacturing and quality, DCMA or production representatives, contractor, supply chain, logistics, cost, and test & EMD before LRIP; repeated in P\&D when production risk changes & Supports LRIP/FRP readiness by checking producibility, manufacturing process control, facilities, tooling, quality, suppliers, and production risk. \\
  OTRR & OTA/COMOPTEVFOR, PM, systems engineer, test, safety, logistics, users, cybersecurity, threat/intelligence, and resource providers & P\&D before OPEVAL/IOT\&E or follow-on OT\&E & Certifies readiness for operational test; for Navy programs, COMOPTEVFOR is the OTRR certification authority for OPEVAL. \\
  PCA & PM, systems engineer, configuration management, contractor, quality, logistics, production, test, and sustainment stakeholders & P\&D, after successful OT\&E and before the FRP or full-deployment decision when tailored & Audits the as-built production-representative system against its technical documentation and supports the final product baseline. \\
  \bottomrule
\end{longtblr}

\subsection{Specifications, Baselines, and Reviews}
\begin{description}
  \item[\textbf{System specifications.}] \ac{srr} reviews whether the system performance specification is clear, testable, and mutually understood; \ac{sfr} approves the system performance specification for use as the functional baseline~\autocite{edo-3-5-3-cm-tech-reviews-2025,DAU-System-Requirements-Review,DAU-System-Functional-Review,DAU-System-Performance-Specification}.
  \item[\textbf{Performance specifications.}] Link to the allocated baseline and \ac{pdr}~\autocite{edo-3-5-3-cm-tech-reviews-2025,DAU-Allocated-Baseline,DAU-Preliminary-Design-Review}.
  \item[\textbf{Detail/process/material specifications.}] Link to the product baseline and \ac{cdr}/\ac{fca}/\ac{pca}~\autocite{edo-3-5-3-cm-tech-reviews-2025,DAU-Product-Baseline,DAU-Critical-Design-Review}.
\end{description}

\subsection{General Principles for Reviews}
\begin{description}
  \item[\textbf{Why.}] Assess design maturity and program risk; not a check-the-box event~\autocite{edo-3-5-3-cm-tech-reviews-2025}.
  \item[\textbf{When.}] Driven by the \ac{sep}, baseline completion, operational test prep, and upcoming milestone reviews~\autocite{edo-3-5-3-cm-tech-reviews-2025}.
  \item[\textbf{Who.}] Led by a \ac{trb} with \ac{pm} representation, engineers, independent reviewers, logistics, cost, legal, contracting, sponsors, and users~\autocite{edo-3-5-3-cm-tech-reviews-2025}.
\end{description}

\subsection{Steps to Effective Design Reviews}
\begin{itemize}
  \item Plan the review and scope~\autocite{edo-3-5-3-cm-tech-reviews-2025}.
  \item Familiarize participants with data and requirements~\autocite{edo-3-5-3-cm-tech-reviews-2025}.
  \item Conduct individual reviews before team review~\autocite{edo-3-5-3-cm-tech-reviews-2025}.
  \item Hold team review and resolve issues~\autocite{edo-3-5-3-cm-tech-reviews-2025}.
  \item Follow up to close actions~\autocite{edo-3-5-3-cm-tech-reviews-2025}.
\end{itemize}

\subsection{Quick Review}
\begin{itemize}
  \item CM avoids cost, reduces risk, and controls design changes~\autocite{edo-3-5-3-cm-tech-reviews-2025}.
  \item CM functions are identification, status accounting, control, and verification/audits~\autocite{edo-3-5-3-cm-tech-reviews-2025}.
  \item Baselines are functional, allocated, and product; changes are engineering changes, deviations, and waivers~\autocite{edo-3-5-3-cm-tech-reviews-2025}.
  \item Commercial items reduce Government control of design details and \ac{tdp} access, increasing interface and support planning needs~\autocite{edo-3-5-3-cm-tech-reviews-2025}.
  \item In \ac{mca}, technical reviews should be explained as risk-reduction gates tied to phase maturity: \ac{srr} reviews requirements in \ac{tmrr}, \ac{sfr} supports the functional baseline, \ac{pdr} supports the allocated baseline before \ac{emd}, \ac{cdr} supports the initial product baseline in \ac{emd}, and \ac{prr}/\ac{otrr}/\ac{pca} support production, operational test, and final product-baseline decisions~\autocite{DoDI5000-88,DoD-EngineeringDefenseSystemsGuidebook-2024,DAU-Technical-Baselines,DAU-System-Requirements-Review,DAU-System-Functional-Review}.
\end{itemize}

%====================
% End of file
%====================
\IfSubfilesClassLoaded{
  \RenewDocumentCommand{\entryname}{}{\textbf{\color{Modern} Acronym}}
  \RenewDocumentCommand{\descriptionname}{}{\textbf{\color{Modern} Definition}}
  \printnoidxglossary[
    type=\acronymtype,
    title=Acronyms,
    style=long-booktabs
  ]}{}
\end{document}
